\section{Discussion}
\label{sec:discussion}

\subsection{Key Findings and Their Implications}

\textbf{Finding 1: Symmetry orbits are geometrically large in real MLP zoos --- the theoretical motivation for canonicalization is empirically confirmed.}

Scaling orbits (mean cosine distance 0.32) and sign-flip orbits (mean cosine distance 1.07) are not geometric curiosities confined to pathological weight configurations --- they are universal properties of the Schürholt MNIST zoo. Every oracle orbit pair we measure exceeds the 0.05 significance threshold. This directly addresses a gap in prior work: \citet{Navon2023EquivariantArchitectures} characterized these symmetries theoretically, but did not measure their empirical size in a real model zoo as cosine distances. Our measurements transform the motivation for canonicalization from a theoretical argument to an empirical baseline.

The implication for practitioners is concrete: any method that processes raw MLP weights for the Schürholt zoo is implicitly contending with a within-orbit diameter of at least 0.32 cosine distance for scaling and 1.07 for sign-flip. Canonicalization --- applied correctly --- removes this noise.

\textbf{Finding 2: NFT is non-invariant to scaling orbits but approximately invariant to sign-flip orbits by construction --- an asymmetry with practical implications.}

The invariance probe reveals that NFT's symmetry handling is asymmetric. For scaling, it is measurably non-invariant (gap$=+0.024$, CI$=[0.0232, 0.0245]$ entirely above zero). For sign-flip, it is approximately invariant by construction (gap$=-0.0007$). This asymmetry was not designed into NFT --- it is an emergent consequence of row-level weight tokenization.

This finding is immediately actionable. For practitioners using NFT-family encoders, scaling canonicalization addresses a measurable capacity waste; sign-flip canonicalization may provide no benefit and, as we show, can be actively harmful when applied with a non-unique algorithm. The recommendation is \textit{scaling-only canonicalization} as the productive preprocessing step for NFT.

The broader implication is methodological: encoder-specific symmetry audits --- probing which symmetries the encoder already handles and which it does not --- should be a standard diagnostic step before designing canonicalization preprocessing.

\textbf{Finding 3: The majority-sign sign-flip canonicalization fails structurally for even-$d_{\mathrm{in}}$ architectures --- a previously uncharacterized limitation.}

The 85.6\% tie rate for $d_{\mathrm{in}}=784$ is a mathematical consequence of $d_{\mathrm{in}}$ being even: with approximately equal positive and negative weights in each row (typical for gradient-trained networks), the expected tie probability per neuron is $\approx 2.8\%$ for $d_{\mathrm{in}}=784$, yielding $\geq 83\%$ probability of at least one tie per 64-neuron model. This structural finding has not appeared in prior weight space learning literature.

\textbf{The fix is clear and implementable.} Three approaches eliminate the tie problem: (1) use an architecture with odd input dimension (e.g., pad MNIST from 784 to 785 --- a single zero-padding that breaks the symmetry with negligible computational cost); (2) use a secondary tie-breaking criterion based on weight magnitude; (3) avoid sign-flip canonicalization entirely for NFT-family encoders, given the emergent sign-flip invariance finding.

\subsection{Limitations}

\textbf{Limitation 1: Zoo scale ($N=500$) provides insufficient statistical power for property prediction experiments.}

All property prediction experiments (H-M2, H-M3) were conducted on $N=500$ models from a local archive. At $n=50$ test samples, Spearman $\rho$ bootstrap CIs have width $\approx 0.6$ --- approximately $12\times$ wider than the effect size we seek to detect ($\Delta\rho \geq 0.05$). No property prediction comparison in this paper has statistical power to distinguish conditions.

This limitation is precisely quantified and does not invalidate our primary contributions. The orbit characterization (H-E1) and invariance probe (H-M1) use relative comparisons that produce tight CIs (width $\approx 0.001$) even at $N=500$. Only the downstream property prediction improvement claim (H-M3) requires large $N$ --- we quantify the requirement: $N \geq 5{,}000$ models.

\textbf{Limitation 2: Sign-flip canonicalization is non-unique for even $d_{\mathrm{in}}$ --- our H-M3 Condition D tested a non-canonical transformation.}

As described in Sections~\ref{sec:results}, the majority-sign algorithm fails for 85.6\% of Schürholt MNIST zoo models. Condition D in H-M3 applied a deterministic but not symmetry-derived transformation to these models. The proper comparison --- clean scaling-only (Condition B) vs.\ Condition E, with adequate $N$ --- has not been performed. We recommend this as the primary follow-up experiment.

\textbf{Limitation 3: The sign-flip functional equivalence in H-M1 used simultaneous two-layer flips; ReLU sign-flip is an approximate symmetry.}

As discussed in Section~\ref{sec:method}, sign-flip functional equivalence for a 2-layer ReLU MLP holds as an approximation. Runtime functional equivalence verification on held-out inputs was not performed and remains a recommended follow-up step.

\textbf{Limitation 4: All findings are derived from a single architecture and dataset; generalizability is unknown.}

The specific numerical findings --- orbit diameters 0.32 and 1.07, NFT gap $+0.024$, tie rate 85.6\% --- are architecture-specific. Whether orbit geometry scales similarly for deeper networks, convolutional architectures, or non-ReLU activations is unknown.

\textbf{Limitation 5: We do not verify whether NFT capacity is the binding constraint.}

At $N=500$, NFT training is severely underpowered, making it impossible to distinguish: (a) NFT is capacity-constrained and canonicalization helps; (b) NFT is underpowered and neither raw nor canonical inputs produce meaningful representations.

\subsection{Future Work}

\textbf{Immediate (directly enabled by this paper):}

\begin{enumerate}
\item \textit{Full Schürholt zoo with scaling-only canonicalization (Condition B).} Acquires adequate statistical power ($N \geq 5{,}000$) and avoids the sign-flip non-uniqueness issue. This single experiment provides a clean test of the core mechanism.

\item \textit{Odd $d_{\mathrm{in}}$ or magnitude-based tie-breaking for sign-flip canonicalization.} Padding MNIST inputs from 784 to 785 eliminates the tie problem and enables a clean test of full canonicalization (Condition D).

\item \textit{Two-layer joint sign-flip functional equivalence verification.} Systematic runtime verification of functional equivalence for the sign-flip oracle pairs used in H-M1.
\end{enumerate}

\textbf{Medium-term:}

\begin{enumerate}
\setcounter{enumi}{3}
\item \textit{Architecture-specific symmetry audit for other encoders.} Apply the same within-orbit vs.\ cross-orbit probing protocol to DWSNets \citep{Navon2023EquivariantArchitectures}, Universal Neural Functionals \citep{Kofinas2024UniversalNeuralFunctionals}, and Hyper-Representations \citep{Schurholt2021HyperRepresentations}.

\item \textit{Deeper networks ($M>2$).} Sign-flip canonicalization for $M>2$-layer MLPs is underdetermined by the majority-sign algorithm. Extending to $M>2$ requires either layer-by-layer greedy canonicalization or a more principled approach.
\end{enumerate}

\subsection{Broader Impact}

This work advances the understanding of geometric structure in neural network weight spaces. The direct applications --- model property prediction, model merging, neural architecture search from weight populations --- are primarily research tools with beneficial applications in model understanding and selection. We identify no significant potential for misuse.

The structural finding on sign-flip canonicalization non-uniqueness is relevant to any weight space learning system that applies sign-flip canonicalization as a preprocessing step. Practitioners should audit whether their architecture has even or odd input dimension before applying the majority-sign algorithm, to avoid introducing structured noise under the guise of canonicalization.
